% Generated draft: requires mathematical and source review.
\documentclass[11pt]{article}
\usepackage[margin=25mm]{geometry}
\usepackage{fontspec}
\setmainfont{Latin Modern Roman}
\usepackage{amsmath,amssymb,mathtools,unicode-math}
\setmathfont{Latin Modern Math}
\usepackage{enumitem,needspace,xurl,hyperref}
\setlength{\parindent}{0pt}
\setlength{\parskip}{5pt}
\setlength{\emergencystretch}{3em}
\newcommand{\N}{\mathbb{N}}
\newcommand{\Z}{\mathbb{Z}}
\newcommand{\Q}{\mathbb{Q}}
\newcommand{\R}{\mathbb{R}}
\newcommand{\C}{\mathbb{C}}
\newcommand{\F}{\mathbb{F}}
\newcommand{\A}{\mathbb{A}}
\newcommand{\PP}{\mathbb{P}}
\newcommand{\E}{\mathbb{E}}
\newcommand{\eps}{\varepsilon}
\newcommand{\dd}{\,\mathrm{d}}
\begin{document}
\section*{UnsolvedMath ID: 30003951}
\textbf{Source problem code:} catalogue entry
\subsection{Definitions and mathematical statement}
\begin{enumerate}
\item A superalgebra is an algebra with a compatible \(\Z/2\Z\)-grading. The first task is to formulate a notion of cluster superalgebra that incorporates spin structures and satisfies a suitable Laurent phenomenon, and to determine whether spectral networks admit a super analogue and what role they play in higher super-Teichmüller theory. In ordinary cluster algebra terminology, a Laurent phenomenon expresses cluster variables as Laurent polynomials in the variables of an initial cluster, with coefficients in a specified coefficient ring. The supplied record does not specify the underlying surfaces, spin structures, coefficient ring, parity conventions, mutation rules, or the proposed super version of this property. It also does not specify the higher super-Teichmüller spaces or the structures that a super spectral network should encode. Thus this is a request for a mathematical framework and construction; the record does not determine a single precise existence assertion.

\item For representations \(\rho_1:H\to\operatorname{GL}(V_1)\) and \(\rho_2:H\to\operatorname{GL}(V_2)\) of the same group \(H\), their direct sum is the representation
\[
(\rho_1\oplus\rho_2)(h)=\begin{pmatrix}\rho_1(h)&0\\0&\rho_2(h)\end{pmatrix}
\quad\text{on }V_1\oplus V_2.
\]
The second question asks whether the direct sum of a quasi-Fuchsian representation and a geometrically infinite representation is stable under small deformations in each of the indicated ambient groups
\[
\operatorname{SO}(6,2),\qquad \operatorname{SO}(4,4),\qquad \operatorname{SL}(8,\R).
\]
Here \(\operatorname{SO}(p,q)\) denotes the determinant-one real linear transformations preserving a nondegenerate symmetric bilinear form of signature \((p,q)\). The related construction problem asks for stable discrete faithful representations that are not Anosov. A representation is faithful if it is injective, and discrete if its image is a discrete subgroup of the ambient Lie group.

The record does not identify \(H\), the two representations, their dimensions, or the embeddings and invariant forms that place their direct sum in the displayed groups. It also does not specify whether the assertion concerns every such pair or a particular pair. A possible meaning of stability here is that the representation lies in the interior of the discrete faithful locus in \(\operatorname{Hom}(H,G)\), with the topology of pointwise convergence; this interpretation is not fixed by the supplied record. The relevant meaning of Anosov, including the parabolic subgroup or family of parabolic subgroups, is likewise unspecified. These missing data prevent an unconditional, fully quantified formulation of this part.

\item Under the usual convention that a finite graph \(\Gamma=(V,E)\) is undirected and simple, its right-angled Artin group is
\[
A(\Gamma)=\left\langle v\ (v\in V)\ \middle|\ [v,w]=1\text{ whenever }\{v,w\}\in E\right\rangle,
\qquad [v,w]=vwv^{-1}w^{-1}.
\]
For \(n\in\Z_{\geq 1}\), let \(\operatorname{SL}_n(\R)\) be the group of real \(n\times n\) matrices of determinant one, and set
\[
\mathcal D(\Gamma)=\left\{n\in\Z_{\geq 1}:\text{there exists an injective group homomorphism }A(\Gamma)\longrightarrow\operatorname{SL}_n(\R)\right\}.
\]
For each finite graph \(\Gamma\), determine the least element of \(\mathcal D(\Gamma)\), including the cases described in the source as unipotent, axial, and pentagon cases. This least-element formulation presupposes that \(\mathcal D(\Gamma)\) is nonempty; no convention for an empty set is supplied. An embedding here means an injective group homomorphism; the record imposes no discrete-image condition on this part.

A matrix \(U\in\operatorname{GL}_n(\R)\) is unipotent if \((U-I_n)^n=0\). The record does not say whether the unipotent case requires the standard generators or all image elements to be unipotent. It gives no definition of axial or the space on which axial behavior is to be measured. The natural interpretation of the pentagon graph is the cycle \(C_5\), with vertex set \(\Z/5\Z\) and edges \(\{i,i+1\}\), but this identification and the intended constraints on its embeddings require source confirmation. The graph convention used above is also an explicit standard interpretation rather than an additional supplied hypothesis.
\end{enumerate}
\subsection{Short English statement}
The record collects three problems: construct cluster and spectral-network frameworks for higher super-Teichmüller theory; investigate deformation stability of specified direct-sum representations and seek stable discrete faithful non-Anosov representations; and determine the smallest real special-linear representation dimension of a right-angled Artin group, including additional special cases. Several essential definitions and hypotheses are absent, so the displayed formulation explicitly retains those uncertainties.
\subsection{Sources}
\begin{itemize}
\item Local UnsolvedMath registry, permanent ID \(30003951\), exact source code \texttt{catalogue entry}, title: \emph{Collected Problems in Higher Teichmüller and Representation Theory}. The supplied record belongs to the Oberwolfach Reports Open Problems collection and gives proposed year \(2018\); no individual proposer is named. Registry locations are \texttt{data/unsolvedmath/problems.json} and the supplementary \texttt{catalog/problems.json}; the supplied evidence does not identify which file contains this entry.
\item Supplied report citation: \emph{New Trends in Teichmüller Theory and Mapping Class Groups} (2019), Oberwolfach Reports, pp.\ 2475--2534, DOI \texttt{10.4171/owr/2018/40}. \url{https://doi.org/10.4171/owr/2018/40}.
\item Supplied literature source, labeled ``Oberwolfach Report 40, Problem 6'': \url{https://publications.mfo.de/bitstream/handle/mfo/3662/OWR_2018_40.pdf?isAllowed=y&sequence=1}.
\item Supplied literature source, labeled ``Kineider et al., Spectral Networks'': \url{https://arxiv.org/abs/2412.03588}. The record supplies no fuller bibliographic details.
\item Literature assessment dated \texttt{2026-08-21}: the supplied curation labels the problems open and states that recent spectral-network work develops related non-super constructions, while the super analogue, stability beyond Anosov representations, and minimal-dimensional right-angled Artin group embedding questions remain open. This is supplied dated curation, not an independently verified literature assessment. The background explicitly requests independent verification of time-sensitive status and cited claims.
\item The supplied registry labels the statement \texttt{corrected\_verified} and reports that the official problem-session report was checked, that three source problems were retained, and that ambient Lie groups and representation-theoretic objects were made explicit. These are supplied verification claims; no source visit or independent mathematical verification is asserted here.
\end{itemize}
\end{document}
